\section{第二层塌陷：真实性危机}

功能层的真空把人推向别处，而别处不是某一个地方，是一组系统。它们的共同点不是提供娱乐，而是提供结果：输入几个词，得到一幅画；发出一个请求，得到一段回应；点击一次，得到一次胜负。得到结果原本需要一个过程，现在这个过程被系统吸收了。变化就从这里开始。

\subsection{低黏滞舒适区：保留刺激，取消做功}

\subhead{完成感的替代品}

留下的问题是：当有用性不再支撑尊严，人会退向哪里。退向的地方不是娱乐，而是一类提供完成感的系统。一个人输入几个词，得到一幅看起来像样的画。他不需要学过透视、构图和色彩，不需要经历几十次画废重来，不需要在某个技术上卡住几个月。他得到的是完成的产物，以及随之而来的、与完成相配的那种反馈。同样的结构出现在写作、作曲、编程、方案策划里。系统输出的是结果，而完成感原本是与过程一起到来的，现在它被单独提取出来，直接交付。

陪伴是同一回事，只是交付的东西更隐蔽。一段持续的关系需要承受对方的情绪波动、需要在自己不想回应的时候回应、需要经历误解和修复。这些成本构成了关系的实际重量，也构成了关系里那种被证明过的确定感。对话系统取消了这些成本，保留了对话中的回应、关注和不被评判的轻松。它交付的是陪伴的感受，而不是陪伴本身。竞技和冒险同样如此。真实的竞争有输的可能，有身体的消耗，有失败后无法撤销的难看。系统里的胜负可以重来，可以调难度，可以在输掉的一瞬间退出。胜利的刺激被保留，随之而来的暴露和代价被削掉。

这三件事表面分散，结构一致：\key{低黏滞系统保留刺激，取消做功；保留感受，削弱后果。}需要说明的是，这个结构不是任何一方设计出来的。没有哪个开发者开会决定要“取消做功”，他们只是各自在优化一个局部指标：提高生成质量、缩短响应时间、提升用户留存。这些优化各自合理，叠加起来却产生了同一个结果，因为被优化掉的恰好都是那些让人付出代价的部分。

\subhead{刺激与后果的解耦}

把上面这个结构说清楚，需要借助一个已有的分析框架，但要先把它从流行的简化版本里拆出来。鲍德里亚用拟像描述过一种状态：符号不再指向现实，而是相互指涉，形成一个自我循环的体系。他的表述常被简化成“假的比真的更像真的”，这个简化丢掉了他论证中最有解释力的部分。\cite{baudrillard1981simulacres}

对本章有用的重构是这样的。拟像的核心不是虚假，而是刺激与后果之间的连接被切断。在真实环境中，一个刺激之所以有意义，是因为它后面跟着后果。火会烫，所以靠近时身体会紧张；承诺会兑现或食言，所以给承诺时人会有分量感；竞争会输，所以投入时人会真正在乎。刺激与后果连在一起，构成一个闭环，闭环让人校准自己的判断。在拟像环境中，刺激被完整保留，后果被移除。感受是真的，手心会出汗，注意力会集中，情绪会起伏，但闭环的另一半不在了。系统给出的反馈与真实反馈在感觉上难以区分，在后果上完全不同。人无法通过这个回路来校准自己：他不会因为判断错了而付出什么，因此也不会更新自己的判断。

高频反馈能否训练判断，要看它在检验什么。反馈若只反映偏好与留存，对外部能力的校准可能很弱；仿真、辅导和游戏也可以通过独立测试产生真实学习。判断标准是迁移、保留与后果，而非频率本身。

\subhead{为什么这不是娱乐批判}

必须把界线划清，否则这个诊断会滑向对娱乐的笼统指责，那既不正确也没有用。娱乐、游戏、虚构、想象，从来都带有脱离直接后果的性质。读小说不为真的经历那些事，看球赛不为自己上场。人类一直在使用低成本的替代性体验来扩展自己的经验范围，这不是新问题，也很难说是一种病。虚构提供了一个安全的试验场，让人在不付出全部代价的前提下接触陌生的处境。

真正的区分不在有没有后果，而在后果链条是否最终闭合。读一整年的小说，人仍然在真实的关系里生活，仍然要面对工作、家庭和身体的后果，虚构是嵌入在闭合链条之中的一段开放区。低黏滞系统的问题不是它提供了开放区，而是它开始占据越来越大的比例，并逐步替代那些原本闭合的环节。当陪伴、成就、竞争和创造都在系统内被满足，真实链条上的对应接口就减少了使用频率，而接口一旦长期不使用，能力会退化，需求也会被重新定向。一个人不再想要真实的关系，不是因为他克制不住，而是因为他已经不再把真实关系识别为满足的来源。

另一个需要排除的解释是道德化的。把这件事说成意志力问题、说成自制力欠缺，会把结构问题转化为个人缺陷，而这个转化正好免除了系统设计者和制度安排的责任。系统之所以有效，恰恰因为它不在道德层面起作用：它提供的满足是真的，成本是低的，反馈是即时的。在这样的条件下，选择低黏滞路径是完全理性的行为。问题不在于个人做了错误选择，而在于当低黏滞路径变得普遍可得时，整个环境的反馈结构发生了改变，而个人的理性适应在这个新结构里产生的结果，和一个有后果的世界里产生的结果不一样。

\subhead{真实感的条件}

如果低黏滞系统的问题是后果被移除，那么真实感的条件就可以被正面地说出来：真实感的来源不是刺激强度，而是刺激与后果之间存在可追踪的耦合。一个判断错了，事情变坏，人因此改变未来的判断，这个过程产生真实感。它不需要剧烈，不需要痛苦，只需要闭合。一个厨师尝了自己做的菜知道咸了，下次少放盐，这就是完整的闭环。闭环的长度可以很短，但两端必须在。

从这个角度看，低黏滞舒适区的危险不在于它让人舒服，而在于它让人舒服而没有被改变。这一区分比它看起来的更难分辨，因为两种体验在主观感受层面可能一致。在一个闭环完整的环境里，舒适的体验也在改变人，因为人从中获得了关于世界的准确信息；在一个闭环缺失的环境里，同一份舒适不留下任何东西，人既不获得技能，也不更新判断，也不建立关系，只是消耗了一段时间。表面上看两种体验都让人愉快，但它们在人的内部留下的东西完全不同。

这就引出一个不易接受的推论：让人感到被满足的环境，不一定是让人成长的环境；而让人感到吃力的环境，也不一定是有价值的环境。判别标准不在感受层面，而在后果是否闭合。这需要把感受和结构分开看，而这两件事在体感上恰恰最难分开，系统也善于生产前者，而前者的信号强度往往压过后者的判断。到这里为止讨论的还是体验的质量，但真正的问题更深一层：如果一个人长期处在后果被移除的环境中，他的判断、他的拒绝能力、他与他人之间的真实联系会发生什么变化，以及他还能不能守住自己的边界。要回答它，得先看清一件事：后果是被什么机制移除的。

\subsection{拟像与反馈断裂}

结构可以被说成刺激与后果的解耦，这是从感受一侧描述的结果。同一件事从系统一侧看，是回路的形状变了：反馈不再由世界发出。要看清这个变化，需要把拟像这个概念从它的流行版本里完整地取出来。

\subhead{符号序列的四个阶段}

鲍德里亚在《拟像与仿真》里给出过一个顺序。符号起初反映深层现实；随后它掩盖并歪曲深层现实；再随后它掩盖深层现实的缺席；到了第四阶段，它与任何现实都不再有关，成为自身的纯粹拟像\cite{baudrillard1981simulacres}。这个顺序不是四种态度，而是连接强度逐级下降的四个阶段，每一级的下降都比上一级更难被察觉，因为符号的表面质量在上升。

这个顺序也不能被读成一张时间表。它不是人类依次经过的四个时代，而是符号与现实之间连接强度的四种状态，同一个社会、同一个领域可以同时处在不同的阶段上。医学期刊与广告文案、工程图纸与竞选宣传，它们与现实的连接强度并不相同；把它们放在同一个阶段里叙述，会得到一个过于整齐的历史图像。这个框架的用处在于逐项判断，而不在于分期。

他用一个寓言开场：一幅地图精确到与领土完全重合，帝国衰败之后，地图的残片留在土地上，人们已经分不清哪一个是真实的\cite{borges1946rigor}。地图先于领土，这就是拟像的判据。判断一个表示环境处在哪个阶段，可以只看一个问题：它能不能被现实打断。能被现实打断的表示仍然是表示；不能被现实打断的表示已经进入第四阶段。

可以把外部支持$g$与表面可信度$q$分开考察。危险情形是$q$上升而$g$下降；这不是所有表示的固定次序，粗糙记录也可能可靠，逼真模型也可能经过充分校准。下式只把需要警惕的一种方向示意出来：

\begin{equation}
g_1>g_2>g_3>g_4\approx 0,\qquad q_1<q_2<q_3<q_4
\end{equation}

$g_k$ 是第 $k$ 阶段符号与现实的连接强度，可以用上段的问题逐项判定，即这个表示能不能被现实打断；$q_k$ 是同一阶段符号的表面质量，包括流畅度、完整度、自洽程度与可信外观。两条序列都是定序量，无量纲，只比较大小，不给出数值。这个式子的含义有两层：连接强度逐级下降，到第四阶段趋近于零，符号只与自身相关；与此同时表面质量逐级上升，两条序列反向。第二层正是“每一级的下降都比上一级更难被察觉”的结构原因：判据在 $g$ 一侧，可观察性却主要在 $q$ 一侧，因此越靠后的阶段越显得正常。适用条件是它只支持逐项判断：同一领域的不同表示可以同时处在不同阶段，两个序列不是时间表，也不能按时代分期；它们刻画的是连接强度的次序，不是某个环境的实测数值。

把这个判据用在内容环境上，可以看得更具体。一条信息是否可被现实打断，取决于它是否连接着一个会付出代价的环节：一条医疗建议如果只需要被转发而无需被验证，它就不进入任何校准；一份评测如果由厂商的预算决定素材，它测的就不是产品。反过来，凡是需要有人付出成本才能产生结果的环节，都自带判决入口：实验要消耗试剂，工程要承担载荷，手术要面对并发症，放贷要对账。可打断性的分布，实际上是代价分布的表征。断裂更容易发生在高流量的表示层：流量奖励的是被接受，而不是被检验。

反馈断裂的机制可以写得更具体。真实世界的反馈有两个属性，替代品很难同时具备：反馈不由接受者生成，反馈不接受协商。一个人做出判断，世界给出的回应与他期待的是否一致，不由他的偏好决定；他也不能要求现实换一个回应再来一次。这两个属性是校准的前提，只有一个信号可能与你期待的不一致、并且不由你控制时，它才携带关于世界的真实信息。在拟像环境里，反馈由系统生成，它的分布可以被调节到与期待一致，不被期待的结果可以被过滤、延后或者改写。校准因此失去发生的条件：\key{要校准，必须先存在一个不服从的项。}

这个机制不限于娱乐环境。任何由被衡量的一方参与生成反馈的系统，都会经历同样的过程：当一项指标被用于决策时，被衡量的一方会调整行为去改善指标，而不是改善指标原本指向的状态。教育评估中的通过率、治安统计中的立案数、金融风控中的评分、科研评价中的论文计数，都出现过这种偏离，因此它以两个名字被记录：经济学的古德哈特定律与社会学的坎贝尔定律\cite{goodhart1975problems}\cite{campbell1979assessing}。这与拟像的四阶段描述的是同一件事，只是发生在组织而不是符号层：信号由被信号描述的一方生产，信号因此不再携带关于状态的信息。

这种断裂不只影响人。系统自身也是从反馈中形成的，模型依赖数据，数据依赖人的活动，人的活动越来越依赖模型的输出。当这一回路足够长，系统就开始从自己的产物中取得材料。前面讨论过这种自我喂养回路在认识论上的后果，这里要看的是它在时间上的代价，因为这一代价可以测量。在递归生成的语料上训练模型，会使模型逐代退化，少数样本携带的尾部信息最先消失，而退化在后续训练中不可逆\cite{shumailov2024ai}。与此并行的是另一项约束：可用于训练的高质量人类文本是有限存量，多项估计认为，这一存量会在本世纪二十年代末到三十年代之间被基本用尽\cite{villalobos2022will}。两件事指向同一个结构：一个主要从自身产物中取得材料的学习系统，会失去与现实接触的比例，而失去接触的代价最终以退化的形式返回。

还有一层更细的断裂发生在测量工具内部。用来评价模型能力的测试集通常取自公开语料，而模型在这些语料上训练过。一个测试问题如果已经出现在训练数据里，它测的就不再是能力，而是记忆；随着版本迭代，基准分数继续上升，分数的含义却在下降。这个问题在机器学习领域被称为数据污染，它看起来是一个纯粹的技术细节，结构与本章讨论的反馈断裂却完全同构：当评价信号的来源与评价对象重叠，信号就失去分辨力\cite{sainz2023nlp}。

\subhead{这个框架的边界}

鲍德里亚关于海湾战争的论述，讨论媒介呈现、战争经验与事件之间的关系。\cite{baudrillard1991gulf} 本书采用的判据是：表示是否保留独立于自身的校验入口。伤亡、破坏与后续行动不会因媒介叙述而消失，它们构成必须进入核算的后果层。

这次校正的代价落在本书这一侧，需要一并说明。要衡量符号与现实的连接强度，前提是现实那一侧可以被独立指认，也就是必须承认存在一个不依赖符号的后果层。这恰好是鲍德里亚否认的命题。本书取走了他的四阶段序列和“下降不再可被察觉”这一诊断，却没有接受让这个诊断成立的那个前提，严格说来用的是他论证的前半段。这个不一致无法靠解释消除，只能被承认：本书的判据（不可逆性、判决入口、后果回流）全部以实在论为前提，如果这个前提不成立，判据随之失效，而不是反过来被论证支持。承认这一点不会让判据变得无用，但它把判据的适用范围限定在愿意接受该前提的讨论之内，而这个范围比“拟像”这个词的流行用法要窄。

这里需要补一个反例，它决定了这个概念该怎么用。符号层与材料层的耦合并非必然断裂，模拟器就是一个正面形态。飞行模拟器、外科训练模拟器和临床试验前的计算机模型，都是低成本的表示环境，它们被允许承担训练与预演的功能，前提是它们被反复地与现实对账：模拟参数来自真实飞行数据，训练效果由真机表现检验，模型预测由试验结果校正。这类系统的关键特征不是更逼真，而是存在固定的对账环节，现实可以随时推翻模拟给出的结论。健康与断裂的分界因此不在表示层的精度，而在有没有一个不由表示层控制的判决入口。

反馈断裂因此有两个受害者。人的判断失去了校准的机会，系统失去了被现实修正的机会。两者的修复方式相同：保留一个必须经过世界的环节。它不能靠意志完成，因为拟像环境的一个特征恰好是不提供拒绝。要看清这一点，需要先回答另一个问题：如果刺激与后果脱钩，欲望的目标是从哪里来的。

\subsection{欲望的镜像与诺斯替复辟}

反馈断裂解释了判断为什么不再被校准，它还没有解释欲望的朝向是怎么被决定的。在低黏滞环境里，人并不是无欲，变化的是欲望的形状，以及形状从哪里来。

\subhead{模仿欲望与中介的距离}

Ren{\'e} Girard 在《浪漫的谎言与小说的真实》里提出过一个结构：欲望不是主体与客体之间的直线，而是一个三角形，主体通过一个中介者去欲求客体\cite{girard1961mensonge}。中介者的存在被一层自我叙述掩盖，人相信自己是在直接欲求某物，而实际上欲求的是他人已经欲求的东西。他把中介分成两类。中介者远离主体、其欲求不构成竞争的时候，是外中介：读一部骑士小说，不会有人与你争夺那位不存在的公主。中介者与主体处在同一个场域、争夺同一个客体的时候，是内中介：竞争随模仿产生，模仿又随竞争强化。

数字环境改变的是中介的距离与数量。算法把中介者从远处搬到近处：你看到的不再是传说中的英雄，而是与你相似的人此刻拥有的东西。中介者的成本同时被压到接近于零，因为展示出来的生活可以省略过程、可以拼接、可以回滚。于是被模仿的是结果的形状，而不是得到结果的过程；模仿的对象越是省略了过程，模仿者的处境与被模仿的结果之间的差距就越大。内中介因此从例外变成默认状态。Girard认为，当所有人的欲望收敛到同一个客体上，冲突无法通过分化解决，只能靠把责任推给一个可以被共同指认的对象来终止\cite{girard1972violence}。这一机制在网络环境中有一个可观察的对应形态：一次集体敌意的形成速度，常常快于事实的核实速度。

中介距离的缩短有一个可以被测量的侧面。社会比较的研究早已指出，人评价自身处境的基准来自周围可比的他人，而不是绝对水平；当比较对象的数量增加、差距被展示得更清楚时，评价基准会整体上移。社交媒体使用与主观感受的研究反复观察到同一个方向：被动浏览他人的展示内容与情绪状态下降相关，而主动的交流并不产生同样的效果。这类研究在因果方向上仍有争议，但结构判断不依赖争议的结论：展示性内容天然省略过程，因此它提供的比较基准是被抬高过的。\cite{festinger1954theory}\cite{verduyn2015passive}

模仿不是病。语言、技能、审美和绝大多数文化传递都依赖模仿，Girard本人也把模仿能力当作人的构成性特征。问题不在模仿，而在中介的接近性与数量结构：中介越远、越少，欲望的空间越宽；中介越近、越多，欲望越收敛到同一组客体上。因此，这一节讨论的是中介的布置，不是人的软弱。

中介的距离、数量、相似性和竞争关系，可以改变人们看到的可能性。更多中介也可能扩展欲望，而非必然使它收敛。Girard的框架提出一条值得检查的机制；具体平台上的集中程度，需要结合内容选择、用户网络与实际行为测量。

\subhead{诺斯替式的逃逸}

当欲望的对象由外部提供，现实这一侧就只剩下代价。走到这一步，会出现一种结构上与古代诺斯替主义同形的判断。诺斯替派是公元二至三世纪的多个教派，它们共享一批主张：世界由一个有缺陷的造物者造成，物质是囚禁灵的容器，得救靠一种关于自身来源的知识，而不靠行动与承受。把它抽象成结构，有三个特征：现实被判为次等；逃离被理解为回归；途径是认识，而不是做功。

这三个特征不需要宗教形式就可以复现。当现实被体验为充满无用的摩擦，而符号层提供低阻力的交互时，把符号层当作更本真的那一层，是一种在结构上划算的判断。它不必宣称现实是假的，只要在实际安排中把现实当作可以绕开的中间层就够了：经验可以生成而不必经历，关系可以模拟而不必承担，作品可以产出而不必训练。这里的关键不是人们相信了什么，而是人们在资源配置上把哪一层当作了主要场所。

这个判据可以写得更具体：一个人在符号层与材料层之间的时间与注意力分配，以及他是否还愿意承担不可回滚的代价。愿意承担的表现形式是有限的几种，长期训练、照护一个具体的人、维护一件会磨损的东西、进入一段无法回退的关系。这些活动的共同点是它们不接受回滚，也不接受由他人代做。判断一个人是否还站在现实这一侧，不需要问他相信什么，只需要看他是否还在做这类事。

本章的论证方向与此相反，理由不在道德，而在结算：\key{物质不是囚禁，而是唯一能承载后果的介质。}符号能生成候选，只有材料、身体和时间能完成结算；代价不是需要摆脱的缺陷，而是位置得以形成的条件。这一点在前面已经出现过：存在的层级越高，维持与回滚的代价越重，而主体恰恰从代价无法转嫁的地方显现出来。

这里需要防止两种误用。一种是把诺斯替复辟当成对虚拟体验的贬低。虚构、游戏和模拟一直是人类扩展经验范围的方式，它们提供的是安全的试验场，不是骗局；问题在比例与闭合，不在于开放区是否存在。另一种是把这一诊断当成道德指控。选择低黏滞路径在给定条件下是理性的：成本低、反馈快、后果可回滚。把它说成意志力问题，恰好免除了环境设计的责任。

当欲望的目标由外部的镜像提供，拒绝的目标也会由外部的镜像提供：镜像崩塌时失去方向，镜像受威胁时产生敌意，镜像需要替罪者时提供替罪者。边界因此不再由内部维持，而由供给决定。走到这一步，问题就从真实层推到了主体层：一个连拒绝的对象都由外部给定的存在，还能不能说自己有边界。
